\chapter{Học Chăm Và Biết Quan Tâm}
\chapterintro{Chủ đề trưởng thành}{Tình cảm ở chương này bớt mơ mộng hơn và gắn với nỗ lực, sự tử tế, chuyện học, chuyện giúp nhau. Vẫn có nét vui nhưng không chỉ còn là nhìn trộm nữa.}

\poem{Bài 41}{Ngũ ngôn tứ tuyệt}{Bảng điểm về cuối tuần}{Em cười khoe điểm tốt}{Tôi chúc nghe rất tỉnh}{Đêm về học nghiêm hơn}
\poem{Bài 42}{Ngũ ngôn tứ tuyệt}{Buổi trực nhật sớm}{Em vẽ mặt cười sương}{Tôi vẽ thêm tim nhỏ}{Xong giả vờ bình thường}
\poem{Bài 43}{Ngũ ngôn tứ tuyệt}{Giờ công nghệ ráp}{Em nối đèn rất xinh}{Tôi nối câu cố gắng}{Vào trong tiếng bạn mình}
\poem{Bài 44}{Ngũ ngôn tứ tuyệt}{Thầy đổi sơ đồ chỗ}{Tôi xa em một hàng}{Mắt vẫn hay đi lạc}{Và ngồi không yên bàn}
\poem{Bài 45}{Ngũ ngôn tứ tuyệt}{Hành lang ôm nắng sọc}{Em ôm chồng sách cao}{Tôi xin cầm bớt nửa}{Vai nhẹ như vừa trao}
\poem{Bài 46}{Ngũ ngôn tứ tuyệt}{Đêm trước ngày kiểm tra}{Em gửi file ghi chú}{Tôi đọc như đọc thư}{Chữ nào cũng biết ủ}
\poem{Bài 47}{Ngũ ngôn tứ tuyệt}{Tiết sử đầy chuyện cũ}{Em chống cằm nghe xa}{Tôi mong mai sau nhớ}{Lớp mình thành chuyện nhà}
\poem{Bài 48}{Ngũ ngôn tứ tuyệt}{Ngày chụp hình cuối khóa}{Em đứng hàng thứ hai}{Tôi đứng rìa mé cỏ}{Mà nắng đậu vào ai}
\poem{Bài 49}{Ngũ ngôn tứ tuyệt}{Con đường me rơi nhẹ}{Em nói mỗi đứa riêng}{Tôi gật như người lớn}{Mà lòng chưa chịu yên}
\poem{Bài 50}{Ngũ ngôn tứ tuyệt}{Khép trang thơ cuối vở}{Tôi viết một câu hiền}{Mai sau dù khác ngả}{Xin còn chào thật hiền}
\poem{Bài 51}{Ngũ ngôn tứ tuyệt}{Nắng leo qua cửa nhỏ}{Em chép bài rất nhanh}{Tôi ngồi ngay phía trước}{Cứ quay về sau hoài}
\poem{Bài 52}{Ngũ ngôn tứ tuyệt}{Mùa thi đi gấp gáp}{Em nhắc nhớ học đi}{Một câu nghe rất nhẹ}{Mà làm tôi bớt lì}
\poem{Bài 53}{Ngũ ngôn tứ tuyệt}{Lưu bút màu mực tím}{Em viết mấy câu ngắn}{Tôi đọc đi đọc lại}{Thành một buổi rất dài}
\poem{Bài 54}{Ngũ ngôn tứ tuyệt}{Giờ địa nhìn bản đồ}{Em hỏi đường ra biển}{Tôi chỉ nhầm sang mắt}{Cả nhóm cười thấy liền}
\poem{Bài 55}{Ngũ ngôn tứ tuyệt}{Cổng trường chen xe sớm}{Em hỏi hôm nay khỏe}{Tôi đáp một câu gọn}{Mà vui nguyên buổi nhé}
\poem{Bài 56}{Ngũ ngôn tứ tuyệt}{Giờ toán tròn tâm nhớ}{Bán kính kéo sang em}{Thầy hỏi đâu là tiếp}{Tôi chỉ ngay môi mềm}
\poem{Bài 57}{Ngũ ngôn tứ tuyệt}{Văn nghệ tập tới tối}{Em giữ nhịp rất đều}{Tôi đàn sai hai nốt}{Em vẫn cười bảo siêu}
\poem{Bài 58}{Ngũ ngôn tứ tuyệt}{Đường sau trường lá rụng}{Em bước chậm hàng cây}{Tôi nghe gió kháo nhỏ}{Tuổi này đẹp biết mấy}
\poem{Bài 59}{Ngũ ngôn tứ tuyệt}{Kiểm tra miệng gọi trúng}{Em nhìn tôi gật đầu}{Tôi trả bài xong thấy}{Cả trời vừa bớt đau}
\poem{Bài 60}{Ngũ ngôn tứ tuyệt}{Cuối kỳ mưa giăng ngõ}{Em trú dưới mái hiên}{Tôi đưa áo khoác mỏng}{Về nhà ấm nguyên đêm}